\chapter{化学平衡不是停住不动}

\begin{kaichang}
桌上放着一瓶没开封的汽水。隔着瓶壁看它，安安静静，一个气泡也没有。你拧开瓶盖——“嘶——”无数气泡争先恐后地往上冲，像一群憋坏了的孩子冲出教室。

请你先猜三个问题的答案，写在纸上：

第一，开瓶之前，那些气体分子是不是“挤在瓶子里憋着，等你开门”？

第二，为什么敞口放上一天，汽水就“跑气”了，变得像糖水一样平淡？气体跑光了吗？

第三，把一瓶汽水放进冰箱，另一瓶放在暖气边，哪一瓶开盖时更“冲”？为什么？

这一章结束时，我们会用一个你多半听过、但极可能理解错了的词来回答它们：\textbf{化学平衡}。很多人以为“平衡”就是停下来、不动了、两边一样多。这一章的任务，就是把这三个“以为”逐个拆掉。
\end{kaichang}

\section{先看见“可逆”这件事}

前面几章我们谈反应，谈的几乎都是“一条道走到黑”：镁带烧完变成氧化镁，不会自己变回来（第 26 章）。但自然界里还有一大类反应是\textbf{双向通行}的：反应物变成生成物的同时，生成物也在变回反应物。这样的反应叫\textbf{可逆反应}。

你其实早就见过可逆过程的迹象，只是没往化学上想：

\begin{itemize}[itemsep=2pt]
\item 汽水、啤酒开瓶冒泡——气体能从水里出来，那它当初是怎么进去的？能进能出，这就是双向。
\item 一杯盐水里盐加到再也溶不动，杯底躺着盐粒——第 25 章管这叫“饱和”。可“溶不动”是真的不再溶解了吗？一会儿你会发现，答案出人意料。
\item 登山的人到高原会喘，回到平原又缓过来——血液抓氧气、放氧气，也是一场双向的拉锯。
\item 冷藏的变色口红、遇热变色的感温杯——颜色变过去还能变回来（不过那是另一类平衡，第 35 章的配位化学会再碰到）。
\end{itemize}

还有一个现象值得专门一提，但你\textbf{只能在视频里看}：一种叫二氧化氮的红棕色气体，封在玻璃管里，放进冰水，颜色明显变浅；放回热水，颜色又变深。颜色来回变，说明管子里有两种物质在互相转化，而且转化方向随温度掉头。二氧化氮有毒，这个实验只能在通风橱里由老师演示，\textbf{绝不要在家尝试}。记住这个画面：温度一变，瓶子里两种东西的“比例”就跟着变——这个比例背后，藏着本章的主角。

\section{把镜头推进去：两个方向的反应同时在跑}

现在把镜头推到分子尺度，看汽水瓶里到底发生了什么。

汽水里溶着大量二氧化碳分子（第 25 章讲过，溶解不是消失，是溶质粒子分散到溶剂粒子中间）。在水面上方、瓶盖下方，有一小段空间，里面也挤满了二氧化碳分子，压强比外面大好几倍——这就是你拧开瓶盖时“嘶”那一声的来源。

关键的画面是：\textbf{水面上，每时每刻都有二氧化碳分子从水里蹦进空气，也同时有分子从空气一头扎进水里}。就像泳池边上，不断有人跳水，也不断有人上岸。没开瓶的时候，跳水的人数和上岸的人数恰好一样多——于是水面（液面上方的气压、水里的浓度）看起来纹丝不动。不是没人动，而是\textbf{进和出打成了平手}。

这就是本章第一个要拆掉误解的地方：看上去的“静止”，在粒子层面是\textbf{动态平衡}——两种方向相反的变化以相同的速度同时进行，宏观上互相抵消，微观上热闹非凡。

用化学的语言说：可逆反应进行到一定程度，正反应（反应物变生成物）的速率和逆反应（生成物变回反应物）的速率变得相等，体系就达到了\textbf{化学平衡}。第 29 章讲过，反应速率取决于有效碰撞的频率；反应物越消耗、浓度越低，正反应就越慢；生成物越积累，逆反应就越快。一慢一快，两条曲线必然在某个时刻相交——那就是平衡点。用一幅人为的示意图来看（坐标轴没有刻度，只表示趋势；曲线形状也是示意，不是实测数据）：

\begin{center}
\begin{tikzpicture}[scale=0.95]
\draw[->] (0,0) -- (7.2,0) node[right] {时间};
\draw[->] (0,0) -- (0,3.6) node[above] {反应速率};
\draw[thick] (0,3.0) .. controls (2,1.9) and (3.5,1.6) .. (6.6,1.55);
\draw[thick] (0,0.15) .. controls (2,1.2) and (3.5,1.45) .. (6.6,1.55);
\node at (1.1,3.25) {正反应};
\node at (1.1,0.55) {逆反应};
\draw[dashed] (3.6,1.55) -- (3.6,0);
\node at (3.6,-0.35) {达到平衡};
\end{tikzpicture}
\end{center}

注意图上一件事：平衡之后两条线都\textbf{没有归零}。它们只是相等了。反应还在跑，永远都在跑，只是两个方向跑的一样快。

\section{速率相等，不等于“一样多”}

这是本章最容易踩的坑，值得单独一节。

“平衡”这个词在日常语言里带着“两边一样重”的暗示——天平平衡，就是两边砝码相等。很多人顺手把这个印象搬进化学：达到平衡，就是反应物和生成物一样多了吧？

完全不是。平衡时相等的是\textbf{两个方向的速率}，不是\textbf{两种物质的多少}。

回到泳池的比方：跳水人数和上岸人数相等，泳池里的人数才稳定。但稳定在一百人还是稳定在十个人，取决于跳水和上岸各自的“难易程度”——池水暖和，大家爱泡在水里，平衡时池里就多；池水冰冷刺骨，大家上岸就不愿下去，平衡时池里就少。两种情况下，进出的“人流量”都相等，池里的“存量”却天差地别。

化学里也一样。有的反应达到平衡时，体系里几乎全是生成物，反应物只剩一点点残渣——我们说它“偏向右边”，看起来就像反应进行到底了；有的反应平衡时几乎全是反应物，生成物只有痕量——看起来就像反应根本没发生。可只要给它时间、把盖子盖严，这两种体系内部都同样在进行着双向的、速率相等的反应。镁燃烧、汽水跑气，在我们眼里是“单向”和“可逆”的区别，在化学家眼里只是\textbf{平衡偏向的程度}差了无数个数量级。

那么，“偏向多少”能不能用一个数字来表达？能。这个数字叫\textbf{平衡常数}。

\begin{qiaoliang}
对一个可逆反应，达到平衡后，把各生成物的浓度相乘、各反应物的浓度相乘，两者相除（每种物质的浓度还要按方程式里的系数取相应次方），得到的商在\textbf{同一温度下总是一个定值}——不管你起始时放的是反应物、生成物，还是随便什么比例，体系最后都会调整到这个数。这个定值就是平衡常数 $K$。

$K$ 的大小直接告诉你平衡偏向哪边：
\begin{itemize}[itemsep=2pt]
\item $K$ 很大（比如 $10^{15}$）：平衡时生成物占绝对上风，反应物几乎消耗干净——反应“几乎进行到底”。
\item $K$ 很小（比如 $10^{-10}$）：平衡时反应物占绝对上风，生成物只有痕量——反应“几乎不发生”。
\item $K$ 在 $1$ 附近：两边势力相当，平衡时谁都吃不掉谁。
\end{itemize}
还有一点必须钉死：$K$ 只随\textbf{温度}改变。改变浓度、压强，只能让体系\textbf{在新的条件下重新凑出同一个 $K$}，而改不动 $K$ 本身。催化剂连碰都碰不到它——第 30 章说过，催化剂只是让到达平衡更快，不改变平衡的位置。
\end{qiaoliang}

\section{扰动来了，体系怎么回应}

现在做本章最实用的思想实验：一个已经达到平衡的体系，你伸手去搅它一下——加点反应物、挤压气体、升个温——它会怎样？

体系不会乖乖受着，也不会完全抵消你的操作。它会\textbf{朝削弱你这个扰动的方向挪动}，直到重新凑回平衡常数要求的那个比例。这个规律叫\textbf{勒夏特列原理}。逐条看：

\textbf{改变浓度。} 向平衡体系里加反应物，正反应这边“原料”突然变多，碰撞变频繁，正反应暂时跑赢了逆反应，于是一部分新加的反应物被转化成生成物，直到两边速率再次打平。反过来，把生成物\textbf{移走}（让它沉淀、流走、被抽走），逆反应就失去了“弹药”，正反应只好继续往前跑，补上缺口。这一条在工业上价值连城，马上会在合成氨的故事里看到。

\textbf{改变压强。} 这一条\textbf{只对气体反应有效}，而且只当反应两边气体分子数目不一样时才起作用。压缩体积，压强升高，体系里气体分子挤得更密、碰撞更频繁；如果某个方向的反应能减少气体分子总数，朝那边走就能“缓解拥挤”。比如氨的合成：左边 4 个气体分子变右边 2 个，加压就把平衡往右边推。要是两边气体分子数相等，加压对两边一视同仁，平衡原地不动。

\textbf{改变温度。} 这是唯一一条会\textbf{改变 $K$ 本身}的扰动。升温相当于给体系添能量：吸热方向的反应像顺风行船，被助长；放热方向则受抑制。于是放热反应升温后 $K$ 变小，吸热反应升温后 $K$ 变大。视频里那管二氧化氮遇冷褪色、遇热加深，就是温度在改写 $K$、拖动平衡的直接证据。

勒夏特列原理用起来快、想得通，是判断方向的好工具。但请你记住它的边界，这条我们放到边界一节细说：它只告诉你\textbf{方向}，从不告诉你\textbf{多少}。

\section{符号登场}

拆到这一层，终于可以写符号了。可逆反应不用单向箭头，而用一对方向相反的半箭头，表示两个方向的反应同时进行。氨的合成就写成：

\[\chem{N_2} + \chem{3H_2} \rightleftharpoons \chem{2NH_3}\]

读法：氮气和氢气变成氨，同时氨也在分解回氮气和氢气；箭头两头通，哪边“车流”大，取决于体系当前离平衡有多远。

二氧化碳进出汽水，写成：

\[\chem{CO_2}\text{（气）} \rightleftharpoons \chem{CO_2}\text{（溶于水）}\]

平衡常数的写法，拿合成氨当例子：生成物浓度的乘积写在分子上，反应物写在分母上，系数变成指数——

\[K=\frac{[\chem{NH_3}]^{2}}{[\chem{N_2}][\chem{H_2}]^{3}}\]

方括号读作“……的浓度”。注意指数：方程式里氨的系数是 2，它的浓度就要平方；氢气的系数是 3，浓度就要立方。这些指数不是凑出来的装饰，它们来自粒子碰撞的计数规则（第 29 章的思路在平衡问题里再次兑现）。纯固体和纯液体（比如汽水瓶里作为主体的水）不写进这个式子——它们的“浓度”在自己看来是个常数，被吸收进 $K$ 里了。

\section{科学家怎样知道：一场赌上粮食与炸药的平衡}

\begin{zhengju}
二十世纪初，欧洲面临一个安静的危机：人口在涨，农田缺氮肥。天然硝石矿快挖完了，而空气里明明飘着装满的氮肥——氮气占空气的约 $78\%$，可惜氮气分子里的三键太牢（第 18 章），植物和人都啃不动它。谁能把空气中的氮“固定”成氨，谁就握住了粮食的钥匙。

有意思的是，勒夏特列本人——就是那位总结平衡移动原理的法国化学家——在 1901 年亲手试过合成氨。他知道按原理应该加压，可实验中混进了空气，高压下氢气与氧气发生爆炸，差点出人命。他就此罢手。

接过这个难题的是德国化学家哈伯。哈伯没有蛮干，他先做的是\textbf{算账}：测出不同温度、压强下这个反应的平衡常数。数字告诉他两个坏消息和一个好消息。坏消息一：常温下 $K$ 小得可怜，平衡几乎完全躺在氮气氢气那边；坏消息二：升温能加快反应，但合成氨是放热反应，升温会把 $K$ 压得更小——速度和位置朝相反方向要温度。好消息是：加压能大幅提高平衡时氨的比例，而且把生成的氨\textbf{冷却移走}、把没反应的氮气氢气\textbf{循环回去}再压一遍，就能用勒夏特列原理里“移走生成物”那一招，逼着反应一轮一轮地往前走。

1909 年，哈伯的装置在约 200 个大气压、约 $500\,^{\circ}\mathrm{C}$ 下稳定地产出氨——温度选在这里，正是“位置要低温、速度要高温”两股要求妥协出来的折中点，再靠催化剂（第 30 章）把速度短板补上。工程师博施随后把它放大成工厂，这就是沿用至今的哈伯—博施法。今天全世界约一半人口的蛋白质里，都有氮原子走过这条管道。

这个故事也要记住它的另一面：同一座合成氨工厂转个产线就能造炸药，它在两次世界大战里都扮演了不光彩的角色。科学原理本身没有立场，方向由人来选。而对我们本章最重要的是方法——哈伯赢在哪里？不是赢在“懂得勒夏特列原理”（勒夏特列本人也懂），而是赢在\textbf{用平衡常数定量地算}：算清楚每个温度压强下平衡到底给他几成氨，再据此设计循环和移走产物的工艺。原理指方向，数字定方案，缺一个都成不了事。
\end{zhengju}

\section{动手试一试：温度怎样摇晃汽水的平衡}

\begin{shiyan}
\textbf{材料：} 同品牌同规格的小瓶汽水两瓶（塑料瓶比玻璃瓶安全）、冰箱、一盆温水（手伸进去觉得温热即可，约 $40\,^{\circ}\mathrm{C}$，绝不烫手）、两个相同的透明杯子、计时器或手机秒表。

\textbf{步骤：}
\begin{enumerate}[itemsep=1pt]
\item 一瓶汽水放冰箱冷藏至少两小时；另一瓶浸入温水中同样时间（敞口放旁边，不要加热密封瓶）。
\item 同时取出，先后快速拧开瓶盖，立刻各倒进一个杯子。
\item 观察并记录：开瓶瞬间“嘶”声的长短和强弱、杯子里气泡涌出的猛烈程度、泡沫持续的时间。
\end{enumerate}

\textbf{观察点：} 哪一瓶开瓶声更响？哪一杯气泡更多、更持久？两瓶在桌上放一小时后，气泡差别是变大了还是变小了？

\textbf{安全提示：} 只许用\textbf{温水}，严禁用开水、微波炉或明火加热密封碳酸饮料——密闭容器内压强骤升可能爆裂。轻拿轻放，不要剧烈摇晃后再开瓶。喝掉实验样品完全没有问题。
\end{shiyan}

你会看到：冷藏的那瓶明显“气更足”——开瓶声更长，气泡更绵密。这正是温度在摇晃平衡：二氧化碳溶进水里的过程是放热的，降温把这个放热过程“往前推”，平衡常数变大，同样一瓶汽水在低温下能溶住更多气体。而温水浴里那瓶，气体早就在瓶内“上岸”挤进了液面上方，水里存不住多少。汽水瓶身上印的“冷藏后饮用更佳”，背后是实打实的平衡常数在变化。

顺手做一个数量级估算，感受一下这个平衡管着多少分子：一瓶 \ud{500}{mL} 的碳酸饮料，出厂时大约压入 \ud{2}{g} 到 \ud{3}{g} 二氧化碳。取 \ud{2}{g} 来算：二氧化碳的摩尔质量约为 \ud{44}{g/mol}，所以物质的量约为 $2\div 44\approx 0.045\ \mathrm{mol}$，乘以每摩尔 \ud{6.02\times 10^{23}}{} 个分子，约等于 $2.7\times 10^{22}$ 个分子。也就是说，你拧开瓶盖的那一刻，有二十多万亿亿个分子正在接受“留下还是出走”的重新裁决——而裁决的标准，就是新的压强和温度下那个 $K$。

\section{身体里也有一套平衡：血红蛋白怎么运氧气}

勒夏特列原理不只管工厂，也管你的呼吸。血液里的血红蛋白是一种能可逆结合氧气的蛋白质（第 15 章提过铁在其中的角色，后面生物大分子部分还会再见面）：氧气分子“搭上”血红蛋白，也能“下车”离开，这是一场标准的可逆过程。

在肺里，肺泡中的氧气分压高——相当于向平衡体系里不断“添加反应物”，于是血红蛋白大量结合氧气，装得满满当当。血液流到肌肉、大脑这些耗氧的组织，那里的氧气被细胞消耗、浓度低——相当于“移走生成物”，平衡掉头，血红蛋白把氧气卸下来，供细胞使用。一装一卸，全靠同一个平衡在两端不同条件下自动换挡，不需要大脑下任何指令。

这也解释了几个生活现象：初到高原，空气稀薄、氧气分压低，血红蛋白装不满，你会喘、会头疼；住上几周，身体增加血红蛋白的制造量，相当于往平衡体系里加“载体”，同样的低氧环境也能运走更多氧，人便慢慢适应。而一氧化碳的危险也在于平衡：它与血红蛋白的结合远比氧气牢固，一旦占据座位就很难让位——这从平衡的角度堵死了氧气的运输线，所以哪怕空气里只有少量一氧化碳也可能致命。家里使用燃气热水器，保持通风不是一句客套话，是在保护你血液里的化学平衡。

\section{这个模型什么时候会失灵}

勒夏特列原理好用，但它是\textbf{判断方向的工具，不是定量计算器}。它能告诉你加压后氨会变多，告诉不了你具体多百分之几——要知道数字，必须像哈伯那样动用平衡常数去算。更麻烦的是，有些情况它会给出\textbf{模糊甚至误导}的答案：多种扰动同时施加时，“哪个方向算削弱扰动”可能说不清；对固体、纯液体参与的平衡，压强条件基本无效；若体系中同时存在好几个相互耦合的平衡（血液里就是这样），单独对某一个用勒夏特列，可能漏掉其他平衡的连锁反应。所以工程上从不靠口诀拍板，靠的是 $K$ 的计算——本章末尾的挑战层会给你一个判断方向的定量办法。

平衡模型本身也有适用范围。第一，它描述的是\textbf{封闭体系}——物质不能进出的那种。敞口的汽水瓶永远到不了真正的平衡，因为气体分子跑掉就不回来了。第二，到达平衡\textbf{需要时间}：有些反应的平衡常数巨大，却慢得几万年也看不出动静（第 29 章），“最终会怎样”和“现在正怎样”是两本账。第三，也是最深刻的：\textbf{生命恰恰不靠平衡活着}。一个达到完全化学平衡的细胞就是死细胞——ATP 消耗殆尽、浓度梯度全部摊平、一切反应两头打平，意味着没有任何净的变化可被用来做功。活细胞维持的是一种不断吃进物质和能量、不断排出废物的\textbf{稳态}：看起来稳定，其实每一分子都在流动，像喷泉的水形而非池塘的水面。这条“远离平衡才有生命”的线索，会在第二册讲代谢时变成主角。

\begin{wuqu}
“反应达到平衡，就是反应停止了，两边物质一样多。”——一句话里两个错误，却是最流行的版本。

反“停止”的证据：把一小块食盐晶体投入一杯\textbf{已经饱和}的食盐水里，按“停止说”，什么都不会发生。但如果水里事先溶有带放射性的钠（科学家用的示踪原子，性质与普通钠几乎一样，只是原子核能发出可探测的信号），过几天检测那块晶体，晶体里出现了放射性。结论只有一个：杯底的晶体在\textbf{不断溶解}，溶液里的离子在\textbf{不断结晶}，两边速率相等，所以晶体的质量不变——但组成它的粒子已经换了一批。平衡是交换的平局，不是交换的终止。

反“一样多”的证据：本章通篇都是。平衡常数大于 $10^{10}$ 的反应，平衡时反应物连亿分之一都剩不下；小于 $10^{-10}$ 的反应，生成物连痕量都难凑齐。它们都可以好好地处在平衡状态——速率相等，数量悬殊。
\end{wuqu}

\section{回到那瓶汽水}

现在兑现开场的三个答案。

第一，开瓶之前气体分子是不是“憋着等你开门”？不是。未开封的瓶子里，溶解与逸出早已打成平手：每秒都有二氧化碳分子冲出水面，也有同样多的分子被压回水里。它们不憋、不等，只是在动态交换中维持着高压下的平衡。你的瓶盖没有“挡住”什么意图，它只是维持了这个平衡存在的边界条件。

第二，敞口放一天为什么跑气？开瓶瞬间，液面上方的高压二氧化碳散入空气，上方浓度骤降，“上岸”这一侧的反应物几乎被清零。平衡被狠狠推了一把，体系只能朝削弱扰动的方向响应：水里的二氧化碳持续不断地逸出，试图重建两相之间该有的比例。可是瓶子敞着口，逸出的分子随风流散、一去不回——这是一个\textbf{开放体系}，真正的平衡永远建不起来，于是单向的净流失一直持续到水里溶解的气体降到与空气相称的极低水平。剩下的糖水不是“平衡了”，是放弃抵抗了。

第三，冰镇的更冲还是温热的更冲？实验已经告诉你：冷藏的那瓶水里锁住的气体更多。因为溶解过程放热，低温把平衡常数推高，同样的瓶内压强下，低温的水能“按住”更多二氧化碳。等你喝到嘴里，体温把汽水焐热，平衡又往回倒，气体在口腔里逸出——那一阵麻酥酥的刺激感，就是几十亿亿个分子集体“上岸”的声音。

\section{继续深挖：水面之下还藏着反应}

汽水的故事其实只讲了一半。溶进水里的二氧化碳并不都以“溶解的气体分子”身份待着，其中一小部分会跟水分子发生真正的化学反应，生成一种能给出质子的物质——碳酸。正是它让汽水尝起来微微发酸，也是它让雨水天然偏弱酸性、让石灰岩在千万年里被溶出溶洞和钟乳石。

“把质子递来递去”的可逆反应，是化学平衡最庞大、最重要的一个家族：它决定了你胃酸的强度、血液为什么不因一口醋就变酸、海洋吸收过量二氧化碳后贝壳为什么会变薄。这一家族有专门的名字和专门的记账尺度，值得我们用整整一章来拆解——那就是第 32 章，酸和碱的故事。

\begin{tiaozhan}
勒夏特列原理的定量版本，其实只需一步：把“当前浓度”而不是“平衡浓度”代入平衡常数的表达式，得到的数叫\textbf{反应商} $Q$。拿合成氨来说，$Q=\dfrac{[\chem{NH_3}]^{2}}{[\chem{N_2}][\chem{H_2}]^{3}}$，用此时此刻的实际浓度计算。

然后只做一次比大小：
\begin{itemize}[itemsep=2pt]
\item $Q<K$：生成物“欠账”，反应向\textbf{正方向}推进，直到 $Q$ 追上 $K$；
\item $Q>K$：生成物“超支”，反应向\textbf{逆方向}退回；
\item $Q=K$：平衡，宏观上不变。
\end{itemize}
勒夏特列说的每一条，都能翻译成这个比较：加反应物，分母变大、$Q$ 变小，于是正向走；移走生成物，分子变小、$Q$ 变小，也是正向走；压缩合成氨的体积，各浓度同倍数放大，但分母的总次方（$1+3=4$）大于分子（$2$），$Q$ 照样变小——加压向右移，原来是指数在起作用。温度则单独行动：它直接改写 $K$ 本身。掌握 $Q$ 与 $K$ 的比大小，你就把勒夏特列原理从口诀升级成了算法。跳过本节不影响主线。
\end{tiaozhan}

\section*{练一练}

\begin{sikaoti}
\item 画一幅“双箱图”：左右两个方框分别代表反应物和生成物，用若干小圆点表示分子，再画两条粗细不同的循环箭头表示正、逆反应的速率。要求画出一种“已达到平衡，但右边分子数远多于左边”的状态，并在图旁注明：这是人为的表示法，圆点大小和颜色不代表分子的真实样子。
\item 一个体系已经达到平衡。某一时刻你突然加入大量反应物。请画一张速率—时间示意图：加入的瞬间，正反应速率和逆反应速率各自怎样跳变？随后两条线怎样重新汇合？（参考本章正文里的图，但要画出“扰动—恢复”的全过程。）
\item 三个反应在同一温度下的平衡常数分别约为 $10^{12}$、$1$、$10^{-8}$。不查任何资料，判断：哪个反应“几乎进行到底”，哪个“几乎不发生”，哪个平衡时两边势均力敌？再回答：给这三个反应都加上催化剂，$K$ 会怎么变？
\item 用平衡移动的语言重新解释：为什么打开的汽水在室温下比在冰箱里“跑气”更快？为什么猛摇汽水后开瓶会喷溅？（提示：摇晃改变了气液接触的机会，先想清楚它动的是速率还是平衡位置。）
\item 初上高原的人呼吸急促，久居后逐渐适应。请画出“氧气—血红蛋白”这一可逆结合的物质流图：在肺、在组织各标出平衡偏向哪边；再解释“身体增加血红蛋白数量”相当于对平衡体系施加了哪种扰动。
\item 一位同学说：“给合成氨的平衡体系加压，平衡右移，说明加压让 $K$ 变大了。”他错在哪里？请分别说清楚：加压改变了什么、没有改变什么；以及哪一类扰动才是唯一能改变 $K$ 的。
\end{sikaoti}
